\section{Classical closures, first jets, and bounded diagram resets}
\label{sec:closure-resets}

Report 28 studies a different way to avoid carrying a large complex through
the entire crossing order.  When a whole relative summand has one matching,
its \emph{actual classical completion} imposes restrictions absent from an
arbitrary algebraic continuation.  In certain cases the entire current
scan has the same total homology rank as that completed matching alone.
We can then splice the matching into the remaining diagram and start a
fresh scan with fewer crossings.  The maintained decision-only backend
\code{closure}, implemented in \code{fastunknot/closure\_scan.py}, combines
these resets with closure-rank obstructions and bounded optional work.
The standard backend remains the default.

\subsection{Source audit and the exact scope of the rerun}

The delivered package distinguishes its shortened reference fixture from
the production checkout.  We ran its AST checker against the current
checkout.  The scanner agrees after docstrings are removed, but the
geometry fixture's \code{compose} omits four maintained fast paths:
zero factors, the two identity cases, and equal endomorphism factors.
The latter uses the square-zero radical identity.  Moreover, the geometry
blob listed in the provenance file is not the blob at the stated inspected
revision \code{13f6dbfdd}: the actual blob there starts \code{92135876},
whereas the supplied pin starts \code{1078526e}.  The other three pins
match that revision.  The earlier intake ledger already flagged this
mismatch, but its summary and catalog wording overstated the verification.
Those summaries are now corrected.  The detailed follow-up is retained in
\code{data/closure-reset-source-audit.json}.

The 29 delivered tests pass with imports redirected to the actual
production package.  We also reran the independent cube audit through
\path{fast/closure_research/audit_upstream.py}.  This adapter changes only
the import paths, output destinations, and scope description; the delivered
files remain unchanged.  All 480 driver comparisons on 160 diagrams,
480 interval-rank checks, and 130 quiver-rank checks pass.  Eighty diagrams
reset and two have early nontriviality certificates.  No nonsingleton
eligible block occurs in this corpus.  The independent cube does not
import scanner code.

\subsection{Why a first jet controls a knot completion}

Let $m$ be a nonempty geometric matching with $p$ arcs, and let a whole
differential component supported on $m$ have coefficient matrix
\[
 Q(\boldsymbol{x})=A+\sum_{i=1}^p x_iB_i+\text{higher dot terms},
 \quad Q^2=0,\qquad B=\sum_iB_i,
\]
over $\F[x_1,\ldots,x_p]/(x_i^2)$.  The matrix raises homological degree
by one.  Gluing the actual suffix gives a link $L_m$, whose Khovanov
complex $C$ carries commuting square-zero dot operators $X_i$.  The
completed block has differential $D=d_C+Q(X)$ on $C\otimes V$.
The component must have no incoming or outgoing attachments; a selected
submatrix does not meet this premise.

\paragraph{The chain-level calculation.}
Moving a marked point across a crossing has reverse-saddle homotopy $H$
with $d_CH+Hd_C=X_a+X_b=\delta$, $H^2=0$, and commutation with all
fixed dot operators.  It also satisfies $H\delta=0$.  These are strict
chain identities, recorded in Nahm's Lemma 2.4 and its cited predecessors
\cite{nahm-basepoints}; equality of induced dot maps on homology would
not suffice here.  For $P=\partial Q/\partial x_i$, differentiation of
$Q^2=0$ gives $PQ+QP=0$.  Thus $G=1+HP(X)$ is self-inverse and
\[
 GDG=D+\delta P(X).
\]
Indeed the linear error is $\delta P$, while the quadratic error is
$Hd_CHP^2=H\delta P^2=0$.  There is no requirement that $P^2=0$.
This conjugation moves one variable's dot action to the new marked edge.

If $L_m$ is a knot, all marks can consequently be moved to one edge.
All higher monomials vanish under $x_i\mapsto X$, leaving $A+XB$.
The characteristic-two vertical operator $\nu$, which changes one circle
label $x$ to $1$ and sums over circles, satisfies
\[
 \nu^2=0,\qquad \nu d_C=d_C\nu,\qquad \nu X+X\nu=1.
\]
Hence $C=E\oplus XE\cong\F[x]/(x^2)\otimes E$, where
$E=\operatorname{im}\nu$ is a subcomplex.  Contracting $(V,A)$ onto
its homology transfers $A+xB$ to $xB_*$: two perturbation factors vanish
because $x^2=0$.  If $t=\dim H(V,A)$ and $b=\operatorname{rank} B_*$,
the transferred complex has dimension $2t$, differential rank $b$,
and homology dimension $2(t-b)$.  Therefore
\[
 \dim H(\operatorname{cl}(Q))
   =\kappa(Q)\,r(L_m),\qquad
 \kappa(Q)=\dim H(V,A)-\operatorname{rank} B_*,
\]
where $r$ denotes total unreduced characteristic-two Khovanov rank.
The formula preserves total rank; it does not assert equality of the
original open complexes or their bigraded profiles.

After exact scalar-unit cancellation, $A=0$, so
\[
 \kappa=N-\sum_h\operatorname{rank} B_h\geq1.
\]
The implementation reads only the parity of singleton-monomial
coefficients in each entry, then computes adjacent binary ranks.  Counting
all monomials modulo two would give the wrong matrix.  Neither barcode
reconstruction nor ranks of products of adjacent matrices are needed.
The maintained reader declines blocks containing scalar units; the general
formula above explains the theory but is not an additional production
scalar eliminator.

\subsection{Pure blocks and the indispensable link distinction}

For a pure block $Q=\theta T$ with one nonzero radical coefficient, put
$\beta=N-\sum_h\operatorname{rank}T_h$ and let
$\sigma=\sum_i[x_i]\theta$ be its total-linear parity.  The scalar
quiver has $\beta$ intervals, although $T^2$ need not be zero.
When $\sigma=1$, the derivation identity
$\nu\theta+\theta\nu=U$ has $U^2=1$ and $\nu U=U\nu$.
Then $U\nu$ splits the suffix complex into a free dual-number factor
for the action of $\theta$.  This proves completed rank $\beta r(L_m)$
for \emph{every} classical link completion.

If $\sigma=0$ and the completion is a knot, the first-jet matrix vanishes,
giving completed rank $Nr(L_m)\geq2N$.  If the completion has at least
two components, each quiver interval contributes at least $r(L_m)\geq4$.
The first inequality follows from the monotone interval profile already
proved in the continuation analysis; the second follows from the link's
Euler specialization.  Consequently a pure block has lower bound
\[
 \begin{cases}
 2\beta,&\sigma=1,\\
 2\min(N,2\beta),&\sigma=0.
 \end{cases}
\]
Any even pure block with an edge therefore certifies rank at least four.
The implementation uses this threshold directly instead of computing
unnecessary exact interval data.

A nonpure block requires the knot-completion premise.  For example, on
the two-component unlink, $x_1+x_2$ acts nontrivially, and a three-term
interval joined by this coefficient has rank four.  Applying the knot
formula would incorrectly give twelve.  The maintained tests include
this counterexample; no completion with several components is silently
treated as a knot.

\subsection{Recognition control, splicing, and fallback}

The backend scans with exact component sharing and multiplicities capped
at three.  At each proper prefix it inspects whole differential components.
Pure components give the closure-uniform lower bounds above.  Otherwise
the driver constructs the actual completion for each eligible matching,
checks its spherical rotation system, and counts link components.
A knot completion supplies $2\kappa$ as a lower bound.  Contributions
are multiplied by their nonnegative integer multiplicities and added;
rank at least four rejects the original knot.

When the entire scan has one component copy, one matching, and
$\kappa=1$ with a knot completion, that residual knot has the same total
rank as the current computation.  Each matching arc joins two boundary
labels of the unprocessed suffix.  The splice merges these labels,
canonicalizes the remaining PD code, and retains the selected crossing
order.  Removed smoothing circles have already been handled by delooping;
they are not reintroduced.  A fresh scanner starts from empty boundary on
the strictly shorter residual diagram.  No claim of a knot isotopy between
the original and residual diagrams is needed or made.

The local \code{closure\_max\_work} allowance, default one million,
charges observer traversal, polynomial extraction, binary elimination,
and completion geometry.  It is an operation allowance, not a bit-cost
or strict wall-clock bound.  The global cooperative deadline has priority.
Local exhaustion disables observation and continues the current exact
scan, preserving all completed algebraic work.  A partially computed
bound or residual is never committed as a reset.  Global exhaustion
returns \code{UNKNOWN} through recognition.  Only final closed rank two
certifies the unknot; a proper-prefix failure to reject is inconclusive.

The decision interface exposes capped ranks and reset events.  It is not
connected to the exact homology or window APIs, because resets discard
grading information.  The optional backend retains the same alternate
backend restrictions as the existing shared scanners.

\subsection{Actual reset gaps and what they do not prove}

Let $g$ be the maximum number of crossings processed between resets or
the final decision, including the initial and terminal segments.  After
$j$ crossings in one segment, there are at most $4j$ frontier labels and,
using a deliberately loose delooping bound, at most $8^j$ physical objects.
The coefficient basis has size $2^{O(j)}$.  Exact arithmetic, elimination,
first-jet inspection, and geometry checks cost
$\operatorname{poly}(n)2^{O(j)}$ in bit operations.  Saturated sharing
does not increase the literal object bound.  The maintained ordering
heuristic has polynomial overhead, and preserving the suffix order means
each original crossing is processed only once across all segments.
Thus this implementation satisfies the parameterized bound
\[
 T(n,g),\ S(n,g)\ \leq\operatorname{poly}(n)2^{O(g)}.
\]
The recorded \code{max\_gap} is this execution parameter, including
scanning after observation has been disabled.  A polylogarithmic bound
on \emph{actual} gaps would give quasi-polynomial recognition.  We do
not know such a bound for arbitrary knots or unknots, or an efficiently
constructible ordering that guarantees it.  Adaptive fallback does not
supply that missing theorem.

\subsection{Expanded discovery, validation, and paired measurements}

The maintained discovery program
\path{fast/closure_research/audit_blocks.py} samples 1200 braid attempts
with seed 2801, accepting 312 knots.  Each is scanned in a random order
and the width heuristic's order, with a quarter-second and 3000-object
allowance per scan.  All 624 scans complete: 4972 crossing stages produce
785 singleton observations and no nonsingleton pure or first-jet block.
This larger negative result supports no claim of actual-diagram gains
from the mixed-coefficient rule.  It does not prove that such blocks
cannot occur.  Inputs, orders, limits and source hashes are retained in
\code{data/closure-reset-block-discovery.json}.

The maintained unit checks directly complete forty mixed polynomial
matrices against the independent trefoil cube, verify $d^2=0$, and compare
their ranks with $\kappa r(L_m)$.  They also compare sixty random knots
against the independent cube under resets, no resets, and zero observation
budget; compare actual splices with the independent implementation;
check malformed blocks and external attachments; and test global deadlines,
CLI dispatch, and local exhaustion after successful resets.  A seven-crossing
Coxeter unknot has six resets and maximum gap one.  This is an illustrative
family already easy for earlier structural recognizers.
The complete maintained suite passes all 567 tests.  Shape caching is tested
both enabled and disabled; automatic mode applies the established repeated-stage
heuristic to each residual diagram, so the raw comparison uses the same cache
policy as the existing saturated scanner.

\begin{center}\small
\begin{tabular}{@{}lrrrrr@{}}
\toprule
Input & Full/reset & Raw/reset & Raw/no reset & Raw/Euler & Raw A/A \\
\midrule
unknot & 0.996 & 0.104 & 0.108 & 0.984 & 0.965 \\
trefoil & 1.028 & 0.953 & 1.126 & 0.991 & 1.180 \\
conway & 0.999 & 0.938 & 0.942 & 0.996 & 1.005 \\
kinoshita\_terasaka & 1.001 & 0.942 & 0.939 & 0.997 & 0.998 \\
hard\_unknot\_8 & 1.022 & 0.889 & 0.923 & 0.961 & 1.001 \\
grid\_scrambled\_unknot & 0.986 & 0.646 & 0.719 & 0.978 & 0.993 \\
stress\_braid5\_36 & 1.001 & 0.927 & 0.949 & 1.014 & 0.990 \\
conway\_sum\_2 & 0.994 & 1.420 & 1.427 & 1.538 & 1.003 \\
conway\_sum\_8 & 1.012 & 5.283 & 5.241 & 5.576 & 1.005 \\
R28 cancel\_pairs\_26 & 1.003 & 0.465 & 0.589 & 0.973 & 1.000 \\
R28 coxeter\_24 & 0.978 & 0.297 & 0.383 & 0.974 & 1.006 \\
R28 figure\_eight & 1.001 & 0.951 & 0.900 & 0.962 & 0.996 \\
R28 positive\_2braid\_11 & 1.015 & 0.940 & 0.941 & 0.993 & 1.004 \\
R28 trefoil & 1.009 & 0.959 & 0.962 & 0.994 & 0.996 \\
R28 trefoil\_then\_curls\_12 & 1.007 & 2.034 & 1.967 & 1.996 & 0.980 \\
R28 trefoil\_then\_curls\_24 & 1.004 & 2.401 & 2.371 & 2.307 & 1.016 \\
R28 trefoil\_then\_curls\_6 & 1.003 & 1.627 & 1.617 & 1.633 & 0.990 \\
torus\_3\_5 & 0.995 & 0.965 & 0.961 & 0.988 & 1.007 \\
torus\_3\_7 & 0.998 & 0.968 & 0.969 & 0.994 & 0.998 \\
torus\_3\_11 & 1.000 & 0.929 & 0.922 & 0.987 & 1.002 \\
\bottomrule
\end{tabular}
\end{center}
Ratios are median paired times; values above one favor the denominator. Full uses the previous default recognizer; raw uses the saturated scanner. The no-reset arm retains closure observations but does not restart.


The full-recognition ratios lie between 0.978 and 1.028, with median
1.001, while the identical full-control ratios span 0.976--1.081.
All twenty inputs are decided by earlier enabled filters, so this corpus
shows no complete-recognition gain from closure dispatch.  The unchanged
default path has median old/new ratio 1.008, with variation consistent
with the controls.  One raw trefoil control reaches 1.180; its small
scanner differences should not be interpreted as reliable gains.

The raw observer improves the Conway sums by about $1.42$ and $5.28$
times, and the trefoil-with-curls examples by $1.63$, $2.03$, and $2.40$
times.  These are early obstructions with \emph{zero resets}; the existing
Euler backend already provides comparable gains.  The actual reset cases
are slower: saturated/reset ratios are 0.889 for the hard unknot, 0.646
for the scrambled grid unknot, 0.465 for cancellation pairs, and 0.297
for the Coxeter example.  Their reset counts are 1, 4, 13, and 23.
The no-reset control is faster on each of these four inputs, showing
that splicing and scanner reinitialization add overhead here.  There is
no observed nonsingleton first-jet compression in this benchmark either.
These measurements justify retaining the optional backend for further
research and leaving automatic recognition dispatch unchanged.

\code{fast/benchmark\_closure.py} measures full recognition separately
from raw saturated, Euler, reset, and no-reset scanners.  It includes the
eight delivered examples and twelve maintained examples, uses seven
shuffled paired rounds, includes validated diagram construction, and
retains an identical control arm in each scope.  Imports and one warmup
are excluded.  Calls have five-second and 50000-object limits; any timed
repetition returning \code{UNKNOWN} suppresses that case's timing ratios.
The old recognition arm loads only \code{recognize.py} from
\code{e6f975521}, sharing unchanged dependencies.  All raw samples and
source hashes are retained in \code{fast/results/closure\_reset\_20261008.json}.
